% !PRINTED 584
\weberpage{584}
%<<CONT
Gleichung (1) kann zwischen den $mn$ Zahlen (2) keine lineare Relation mit rationalen Zahlencoëfficienten bestehen.

Jede Zahl $\omega$ des K\"orpers $\Omega$ kann als ganze Function $(n-1)^{\text{ten}}$ Grades von $\Theta$ mit Coëfficienten in $R$ dargestellt werden, und jede solche Zahl $\omega$ gen\"ugt einer in $R$ irreduciblen Gleichung h\"ochstens vom $n^{\text{ten}}$ Grade; $\omega$ ist dann und nur dann eine ganze Zahl, wenn diese Gleichung, nachdem der Coëfficient der h\"ochsten Potenz auf 1 gebracht ist, ganze Zahlen in $R$ zu Coëfficienten hat.

Den Inbegriff aller ganzen Zahlen in $\Omega$ bezeichnen wir, wie fr\"uher, mit $\mathfrak{o}$.

Wenn wir, ohne die Zahlen des K\"orpers $R$ zu ver\"andern, f\"ur $\Theta$ die s\"ammtlichen Wurzeln $\Theta_1, \Theta_2, \ldots, \Theta_n$ der Gleichung (1) nehmen, so erhalten wir $n$ K\"orper
\weq{3}{\Omega_1 = R(\Theta_1), \ \Omega_2 = R(\Theta_2), \ \ldots, \ \Omega_n = R(\Theta_n),}
\gesperrt{die in Bezug auf den K\"orper $R$ conjugirt sind.}

Sind $\omega_1, \omega_2, \ldots, \omega_n$ entsprechende Zahlen dieser K\"orper, so heisst das Product
\weq{4}{\mathfrak{N}_R(\omega) = \omega_1 \omega_2 \ldots \omega_n}
\gesperrt{die in Bezug auf $R$ genommene Partialnorm (Relativnorm)} der Zahl $\omega$.

Eine Zahl $\eta$ des K\"orpers $R$ hat eine in diesem K\"orper genommene gew\"ohnliche Norm $N_R(\eta)$, die eine rationale Zahl ist. Die Relativnorm $\mathfrak{N}_R(\omega)$ ist nun eine solche Zahl $\eta$, und wenn wir ihre Norm im K\"orper $R$ nehmen, so erhalten wir die \gesperrt{Totalnorm} der Zahl $\omega$ im K\"orper der rationalen Zahlen:
\weq{5}{N_\Omega(\omega) = N_R\,\mathfrak{N}_R(\omega).}

Dies ergiebt sich, wenn man $\omega$ durch die Potenzen (2) von $\Theta$ und $\varrho$ ausdr\"uckt, sodann f\"ur $\varrho$ die $m$ conjugirten Werthe, und zu jedem dieser $\varrho$ die nach $R$ conjugirten Werthe $\Theta$ setzt.

In demselben Sinne haben nun auch die \gesperrt{Functionale} des K\"orpers $\Omega$ und die durch diese repr\"asentirten \gesperrt{Ideale ihre Partialnormen}. Die Partialnorm eines Functionals in $\Omega$ ist ein Functional in $R$, und demnach ist die Partialnorm eines Ideals in $\Omega$ ein Ideal in $R$.

Von den conjugirten Idealen k\"onnen auch zwei oder mehrere einander gleich sein, und zwar kann dies auch dann eintreten, wenn die die Ideale repr\"asentirenden conjugirten Functionale nicht identisch sind, wenn sie nur associirt sind. Wenn man
%CONT>>
